% =====================================================================
%  Professioni 4 — Il rigger, l'animatore e il technical animator
%  Fonte: trascrizione della presentazione del corso di character
%  animation ai Dev Days (AIV4): un incontro online di circa un'ora,
%  in tarda mattinata, tenuto da Yuri Perrini, coordinatore del corso di
%  Character Animation di AIV. Circa 22 partecipanti, in parte del primo
%  anno e in parte di anni successivi; slide e siti condivisi a schermo;
%  ampia parte finale di domande, a voce e in chat.
%  Pacchetti richiesti nel preambolo del libro:
%    \usepackage[utf8]{inputenc}, \usepackage{booktabs},
%    \usepackage{tikz} con
%    \usetikzlibrary{arrows.meta,positioning,calc,fit}
%
%  Collocazione nel libro: parte "Alcune Professioni nel Mondo
%  Videoludico", dopo AIV3_narrative_design (3.5).
%  Tutte le etichette di questo capitolo hanno il prefisso "p4:".
% =====================================================================

\chapter[Rigger, animatore e technical animator]{Il rigger, l'animatore e
il technical animator: la character animation per i videogiochi}
\label{p4:cap:animation}

Fra le presentazioni di corso dei Dev Days (capitolo~\ref{p1:cap:devdays})
c'è quella di \textbf{Character Animation}, tenuta dal coordinatore del corso,
\textbf{Yuri Perrini}. Come le altre, non serve tanto a spiegare il corso ai
suoi iscritti quanto a far capire a tutti \textbf{come collaborare con
animatori e rigger} durante i progetti interdisciplinari (i Dev Projects,
sezione~\ref{p1:sec:devprojects}), quali figure professionali possono uscire
dal corso e ``qual è l'anima'' di ciascuna specialità.

L'incontro si tiene online. All'inizio il relatore conta ``una ventina'' di
partecipanti, poi ``intorno ai 20, intorno a 22, compreso me''. Non sono
tutti del primo anno: alcuni vengono da una vecchia scuola chiamata ``Event
Horizon''\footnote{Così nella trascrizione; il nome non è verificabile. Dal
contesto, una scuola confluita in AIV o i cui studenti sono stati
ricollocati in AIV ad anni successivi al primo.} e sono al secondo anno,
alcuni programmatori addirittura al terzo. La trascrizione non riporta né la
data né l'ora; il relatore chiude dicendo che i ragazzi ``hanno fame e
vogliono andare a pranzo'', quindi l'incontro si tiene in tarda mattinata,
con ogni probabilità in uno dei giorni della settimana dei Dev Days, fra il
5 e il 9 ottobre 2026\footnote{Per la deduzione delle date della settimana si
veda il capitolo~\ref{p1:cap:devdays}. Anche il relatore conferma l'anno: dice
di lavorare nell'animazione ``dal 2009 circa, quindi sono 17 anni''.}.

Il relatore avverte subito che vorrebbe una lezione ``molto, molto
interattiva'': visto che il pubblico non è fatto di studenti di character
animation, invece di un \emph{deep dive} sull'animazione e sul rigging,
``che magari a voi interessano parzialmente'', preferisce che gli si chieda
come quelle figure possano essere utili: ``se devo progettare questo
videogioco, cosa mi serve? Quanto può produrre un animatore, per sapere
quanti animatori mi servono per il mio progetto?''. E infatti la seconda
metà dell'incontro è dedicata all'ingresso nell'industria (portfolio,
curriculum, stipendi, test) e alle domande.

Avvertenza di lettura: la trascrizione è automatica. Si apre con una frase
allucinata dal software (``Signora Presidente, onorevoli colleghi'') e
contiene almeno due blocchi di testo senza senso ripetuti più volte (uno a
metà del discorso sulle figure professionali, ``Quindi, personalmente,
questa è la prima cosa che succede'', e uno verso la fine, ``la capacità di
impostare la loro dottrina sul mercato del lavoro''), che qui sono stati
ignorati. Storpia inoltre molti nomi: AIV diventa ``Ive'', ``Dive'',
``Hive''; ArtStation ``RStation''; \emph{The Last of Us} ``delasto vase'';
Unreal Engine ``Arial Engine''; il technical animator ``Technicolimator''.
Dove la ricostruzione è sicura si usa la forma corretta; dove non lo è lo si
dice in nota. Durante la parte sull'industria il relatore mostra a schermo
alcuni siti (un foglio Google con le offerte di lavoro, un foglio con gli
stipendi, Numbeo) e i profili LinkedIn di alcuni conoscenti: le cifre
riportate sono quelle lette a voce in quel momento.

Il percorso del capitolo:
\begin{enumerate}
  \item il corso e il relatore (sezione~\ref{p4:sec:corso});
  \item lo \textbf{spettro} delle figure, dall'artistico al tecnico, e il
        technical animator (sezione~\ref{p4:sec:spettro});
  \item dove si colloca il rigging e l'animazione nella \textbf{pipeline}
        degli asset, e il lavoro in parallelo
        (sezione~\ref{p4:sec:pipeline});
  \item cambiare specializzazione si può (sezione~\ref{p4:sec:cambiare});
  \item dal rig all'animazione: gameplay e cinematiche
        (sezione~\ref{p4:sec:animazioni});
  \item il \textbf{programma} dei due anni (sezione~\ref{p4:sec:programma});
  \item entrare nell'industria: portfolio, network, curriculum, cover
        letter (sezione~\ref{p4:sec:industria});
  \item le offerte di lavoro, gli stipendi, il costo della vita
        (sezione~\ref{p4:sec:stipendi});
  \item l'età, la provenienza, gli pseudonimi, il portfolio
        (sezione~\ref{p4:sec:eta});
  \item genere e negoziazione (sezione~\ref{p4:sec:genere});
  \item i \textbf{test} di ammissione (sezione~\ref{p4:sec:test});
  \item le domande finali (sezione~\ref{p4:sec:domande}).
\end{enumerate}

% ---------------------------------------------------------------------
\section{Il corso e il relatore}
\label{p4:sec:corso}

\subsection{Un corso verticale}

Character Animation è ``uno dei corsi un po' più giovani'' di AIV, creato più
di recente per rispondere a un'esigenza precisa: formare figure
\textbf{specializzate in maniera verticale sul rigging e sull'animazione per
i videogiochi}. Per due anni gli studenti si concentrano principalmente su
rigging e animazione, con l'eccezione di qualche modulo propedeutico, pensato
per capire il workflow e la pipeline intera dello sviluppo:
\begin{itemize}
  \item un modulo iniziale di \textbf{modellazione}, ``solo un po' per
        prendere dimestichezza col software'';
  \item un modulo di \textbf{Unity} e uno di \textbf{Unreal}, ``senza
        ovviamente avere pretese di rubare il lavoro ai programmatori, ai
        level designer, ai game designer'', ma per capire come integrare le
        proprie animazioni e i propri rig nel motore.
\end{itemize}

\subsection{Chi è il relatore}

Perrini lavora nell'animazione \textbf{dal 2009 circa}, ``quindi sono 17
anni''. Ha vissuto in \textbf{Islanda e in Canada}, lavorando a produzioni
come \emph{Stranger Things}, \emph{Il gatto con gli stivali} e \emph{Rebel
Moon}\footnote{Così nella trascrizione; il relatore non specifica a quale
titolo della serie del Gatto con gli stivali, né quale sia stato il suo
ruolo nelle singole produzioni.}: principalmente televisione e cinema. Ha
fatto anche varie esperienze nel videogioco:
\begin{itemize}
  \item con \textbf{Stan Lee}, che prima di morire\footnote{Stan Lee è morto
        il 12 novembre 2018.} ha creato \textbf{Stan Lee's Kids Universe},
        ``una sorta di app con vari videogiochi dentro per ragazzi'' sulle
        sue proprietà intellettuali, ``non solo Marvel'';
  \item con i ragazzi di un piccolo studio\footnote{Nella trascrizione
        ``Ravig Studio'': il nome non è verificabile, come quello del gioco
        (``Fade Out'').} che stanno per lanciare su Steam \emph{Fade Out},
        un horror cooperativo multigiocatore indie\footnote{Nella
        trascrizione ``un doppia indie horror co-op multiplayer'': forse
        ``un double-A indie'', oppure ``un horror indie co-op''. Il
        relatore ne mostra il trailer e mette il link in chat.};
  \item con \textbf{One-O-One Games}\footnote{Nella trascrizione ``101
        Games''. È lo studio milanese di \emph{Il suicidio di Rachel
        Foster} (2020), citato più avanti.}.
\end{itemize}
Online si presenta soprattutto con la parte TV e cinema (``se mi cercate
online, probabilmente trovate anche dove ho fatto le scuole elementari''),
quindi le esperienze nel game non si trovano facilmente. Ha fondato uno
studio\footnote{Nella trascrizione ``Mostry for You''; la grafia esatta non è
verificabile.} con cui sviluppa prodotti \emph{real time}, non solo
videogiochi ma anche AR e VR, e offre servizi di assistenza alle produzioni
di serie TV, cortometraggi, pubblicità.

Insegna anch'egli dal 2009: ha cominciato a lavorare e a insegnare
contemporaneamente, ``in una miriade di posti''. Da 12 anni è
all'\textbf{IED di Roma}\footnote{Nella trascrizione ``all'OIED di Roma'':
l'Istituto Europeo di Design.}, dove è anche coordinatore; ha insegnato
all'\textbf{American University of Rome}; in una scuola online ``paragonabile
o simile ad AIV'', ma ``di base australiana''\footnote{Il nome della scuola
cade nella trascrizione. Una scuola online australiana di computer grafica
molto nota è CG Spectrum, ma l'identificazione è solo un'ipotesi.}; e ad
\textbf{Animation Mentor}, scuola online di animazione che ``forse qualcuno
di voi conoscerà''. Insegnare è una cosa che fa ``molto volentieri per
passione'', come piccola parte della settimana lavorativa: di solito non più
di \textbf{15 ore alla settimana}.

\subsection{Le due anime del corso}

Lo studio l'ha fondato con il socio \textbf{Daniele Dolci}, che è anche
co-ideatore del corso di Character Animation insieme ad Andrea
Morini\footnote{Andrea Morini è uno dei due relatori dell'incontro di
apertura dei Dev Days (capitolo~\ref{p1:cap:devdays}).}: gli iscritti al
corso trovano entrambi. Il relatore lo racconta perché le due persone
rappresentano bene \textbf{le due anime del corso}:
\begin{itemize}
  \item Perrini, nonostante uno spiccato senso tecnico (``probabilmente sono
        un programmatore mancato'': da giovanissimo faceva il web developer),
        ha sempre avuto un'anima più \textbf{creativa e artistica}. Si è
        avvicinato al 3D pensando di fare scultura e modellazione di
        character e creature; poi ha scoperto l'amore per l'animazione; ma
        sia lo sculpting sia l'animazione erano, nella sua testa, ``servi''
        dell'idea di fare la \textbf{regia} di un prodotto. Negli ultimi anni
        è diventato regista: sta per uscire una serie animata che ha
        interamente diretto, \emph{Melody e Momon}\footnote{Così nella
        trascrizione; il titolo non è verificabile.}, 52 episodi da 7
        minuti, sulla Rai;
  \item Dolci è un \textbf{rigger}, figura più tecnica: ingegnere informatico
        e programmatore, con una laurea triennale in ingegneria informatica e
        una magistrale\footnote{Nella trascrizione ``una laurea magistrale in
        matematica astratta, ingegneria avanzata, AI'': l'elenco è confuso e
        non permette di ricostruire con certezza il titolo del corso di
        laurea.} con tesi su machine learning e intelligenza artificiale
        applicati alla computer grafica.
\end{itemize}
Due persone, dunque: una con un'anima \textbf{tecnico-scientifica}, una con
un'anima \textbf{creativo-artistica}.

% ---------------------------------------------------------------------
\section{Dall'artistico al tecnico: lo spettro delle figure}
\label{p4:sec:spettro}

Le figure della produzione di un videogioco si possono disporre lungo uno
\textbf{spettro} (figura~\ref{p4:fig:spettro}):
\begin{itemize}
  \item a un estremo, le figure più artistiche: \textbf{modellatori,
        texture artist, concept designer}, persone che vengono magari dal
        disegno o dalla scultura;
  \item poi gli \textbf{animatori}, che oggi, soprattutto nel videogioco,
        sono persone appassionate alla parte creativa e artistica ma che
        devono comunque conoscere una parte tecnica;
  \item poi il \textbf{technical animator}, figura che esce dal corso;
  \item all'altro estremo, le figure tecniche ``vere e proprie'': il
        \textbf{rigger}, il programmatore, il level designer, chi fa
        Blueprint, chi imposta le \emph{state machine}.
\end{itemize}

\begin{figure}[htbp]
\centering
\begin{tikzpicture}[>=Stealth, font=\small,
  box/.style={draw, rounded corners, align=center, minimum height=1cm,
              text width=2.2cm, font=\scriptsize}]
  \draw[<->, thick] (0,0) -- (12,0);
  \node[below, font=\scriptsize] at (0.6,-0.1) {creativo-artistico};
  \node[below, font=\scriptsize] at (11.4,-0.1) {tecnico};
  \node[box, anchor=south] (a) at (1.3,0.55) {concept designer,\\modellatori,\\texture artist};
  \node[box, anchor=south] (b) at (4.4,0.55) {animatori\\(creativi, ma con\\base tecnica)};
  \node[box, fill=gray!15, anchor=south] (c) at (7.5,0.55) {technical\\animator\\(ibrido)};
  \node[box, anchor=south] (d) at (10.6,0.55) {rigger, programmatori,\\level designer,\\Blueprint, state machine};
  \foreach \n/\x in {a/1.3,b/4.4,c/7.5,d/10.6} \draw (\x,0.55) -- (\x,0);
\end{tikzpicture}
\caption{Lo spettro delle figure descritto dal relatore, dall'anima
creativo-artistica a quella tecnica; il technical animator sta nel mezzo.
Ricostruzione schematica.}
\label{p4:fig:spettro}
\end{figure}

\subsection{Il technical animator}

Il relatore chiede chi sa che cos'è un \textbf{technical animator}, o chi non
ne ha mai sentito parlare. È una figura che esiste \textbf{principalmente
nel videogioco}: qualcuno che conosce un po' la parte tecnica (in Unreal
Engine, in Unity), un po' quella artistica, e anche la parte di
\emph{effects} e simulazioni. In una produzione il technical animator è la
persona che:
\begin{itemize}
  \item \textbf{aggiusta le cose} e si assicura che tutto venga esportato in
        maniera corretta;
  \item fa le \textbf{simulazioni}: per esempio, in \emph{The Witcher}, i
        capelli lunghi, i vestiti, le code, ``cose che non sono animate a
        mano'';
  \item si occupa di \textbf{retargeting}, cioè di trasferire il
        \emph{motion capture} su personaggi \emph{custom}.
\end{itemize}
È una figura ibrida fra rig, animazione, pipeline e un po' di simulazione:
un mix di creativo, artistico e tecnico.

% ---------------------------------------------------------------------
\section{Dove si collocano rigging e animazione nella pipeline}
\label{p4:sec:pipeline}

Il relatore mostra una schermata con i corsi di AIV disposti come una
pipeline. Le \textbf{animazioni e i rig vanno pensati come asset}, al pari
dei modelli. Il percorso di un asset è (figura~\ref{p4:fig:pipeline}):
\begin{enumerate}
  \item si parte dal \textbf{concept art} (che è anche un corso): il design
        illustrato dei modelli;
  \item si passa per \textbf{videogame 3D art}, cioè modellatori, texture
        artist, shading artist, lighting ecc., che rendono tridimensionale e
        utilizzabile un concept che di solito è solo 2D (si veda il
        capitolo~\ref{p6:cap:3d});
  \item già da qui \textbf{comincia una discussione con il rigger},
        principalmente, ma anche un po' con l'animatore, sull'anatomia del
        personaggio, su dove devono stare le ossa, su come il rigger vuole la
        topologia;
  \item infine si passa a \textbf{character animation}.
\end{enumerate}

\begin{figure}[htbp]
\centering
\begin{tikzpicture}[>=Stealth, font=\small, node distance=0.6cm,
  box/.style={draw, rounded corners, align=center, minimum height=0.9cm,
              text width=2.4cm, font=\scriptsize}]
  \node[box] (c) {concept art\\(design 2D)};
  \node[box, right=of c] (m) {videogame 3D art\\(modello, texture,\\shading, luci)};
  \node[box, right=of m] (r) {rigging\\(ossa, controlli,\\blend shape)};
  \node[box, right=of r] (a) {animazione\\(gameplay e\\cinematiche)};
  \draw[->] (c) -- (m); \draw[->] (m) -- (r); \draw[->] (r) -- (a);
  \draw[<->, dashed] (m.south) .. controls +(0,-0.8) and +(0,-0.8) ..
    node[below, font=\scriptsize]{discussione su anatomia, ossa, topologia} (r.south);
  \node[box, text width=10.6cm, above=1.1cm of $(m)!0.5!(r)$] (p)
    {in parallelo: storytelling/narrative design (dal concept), game design e
     programmazione (con asset \emph{proxy}, poi sostituiti con quelli
     definitivi)};
  \draw[->, dotted] (p.south -| c.north) -- (c.north);
  \draw[->, dotted] (p.south -| a.north) -- (a.north);
\end{tikzpicture}
\caption{La pipeline degli asset come la descrive il relatore: le animazioni e
i rig sono asset come gli altri. Ricostruzione schematica della schermata
mostrata.}
\label{p4:fig:pipeline}
\end{figure}

\subsection{La topologia e le preferenze del rigger}

Il relatore chiede chi è di videogame 3D art: rispondono Giulia, Stefano,
Alessandro; e chi ha già un'idea della specializzazione che preferisce.
Giulia dice \emph{character and environment}. A lei, e a tutti i futuri
modellatori, il relatore dà un avvertimento: nel corso verranno insegnate
delle \textbf{best practice sulla topologia}; ``preparati: verranno
completamente ribaltate da ogni rigger che incontrerai nella tua carriera''.
Ogni rigger ha le sue preferenze: ``no, non voglio quell'edge loop a 5 punte
là, ma ne voglio una a 3 più qua''\footnote{Il riferimento è ai
\emph{poli} della topologia, cioè ai vertici in cui convergono tre o cinque
spigoli invece dei quattro di una griglia regolare; la loro posizione
influisce su come la mesh si deforma.}. Bisogna imparare a capire
\textbf{che cosa} si sta chiedendo e \textbf{perché}: la topologia di un
modello definisce come il modello si deforma, e a seconda di come il rigger
imposta il setup ha bisogno di cose diverse per evitare che si deformi troppo
o troppo poco.

\subsection{Il lavoro in parallelo}

Il passaggio da 3D art a character animation è uno dei principali momenti in
cui \textbf{si comincia a lavorare in parallelo su più fronti}:
\begin{itemize}
  \item la scrittura\footnote{Nella trascrizione ``strategic writing'', e
        poco dopo ``Narrative Design fa parte, nell'ambito videogioco, di
        Strategic Writing''. Con ogni probabilità è la storpiatura del nome
        di un'area o di un corso di scrittura (\emph{storytelling and
        writing}?), di cui il narrative design sarebbe la declinazione per i
        videogiochi; non è verificabile.} sarà già partita insieme al
        concept; il narrative design ne è la declinazione per il videogioco;
  \item game programming e game design saranno partiti più o meno in
        contemporanea allo sviluppo degli asset. Di solito si comincia a
        sviluppare game design e programmazione con degli \textbf{asset
        proxy}: non il personaggio definitivo, non il rig definitivo, non
        l'animazione definitiva, che vengono poi sostituiti (\emph{swap}) con
        quelli giusti.
\end{itemize}
Da un certo punto in poi è come se tutti i corsi lavorassero in
contemporanea. Immaginando di essere in Ubisoft, ognuno dei corsi rappresenta
uno step produttivo, \textbf{un dipartimento} dello studio.

% ---------------------------------------------------------------------
\section{Cambiare specializzazione si può}
\label{p4:sec:cambiare}

Nel mondo del lavoro, sviluppando competenze oltre ciò che si è studiato, ci
si può occupare di più cose, soprattutto negli studi di dimensioni medie e
medio-piccole (``dalle 150 persone in giù''). Anzi, \textbf{sapersi occupare
di più cose rende più difficilmente sostituibili}.

Due esempi:
\begin{itemize}
  \item Andrea Priore, un buon amico conosciuto quando lavorava in
        \textbf{Reply}, dove è stato 13 anni, occupandosi principalmente di
        modellazione e di arte 2D e 3D; si è poi dato alla character
        animation ed è diventato animatore ``all'interno della stessa
        azienda''; ora è in Ubisoft\footnote{La trascrizione è contraddittoria
        (``adesso sta Ubisoft'', poi ``adesso è diventato animatore in
        Ubisoft all'interno della stessa azienda''): non è chiaro se il
        passaggio da modellatore ad animatore sia avvenuto in Reply o in
        Ubisoft. Il relatore mostra il profilo LinkedIn.};
  \item un suo insegnante, David Weatherly\footnote{Grafia del cognome
        incerta.}, è stato in \textbf{DreamWorks} ``per una vita'': faceva
        inizialmente i rig, poi è diventato animatore attraverso una sorta di
        test, di concorso interno.
\end{itemize}

Da studenti si è portati a fare una scelta, ma anche se ora si è scelta una
strada \textbf{non è impossibile cambiarla}. Per ragioni esterne (``arriva
l'AI, la modellazione cala tantissimo, e allora io che mi sono dato alla
modellazione decido di darmi all'animazione, che invece vedo che cala di
meno'': ``ipotesi del tutto campata in aria'') o perché ci si annoia, come è
successo al relatore. Lui si è annoiato della modellazione: ``una volta che
ho modellato un umano, un bambino, un cane, un gatto, sembra che sto rifacendo
versioni diverse della stessa cosa''. Come animatore, invece, ha
l'impressione di essere un \textbf{attore}, di recitare ogni volta una parte
diversa, soprattutto nell'animazione per le cinematiche, dove il personaggio
deve parlare e mostrare emozioni (``immaginatevi un \emph{The Last of
Us}''): si fa comunque una performance, soprattutto quando il personaggio è
cartoon, e quindi \textbf{senza motion capture}.

% ---------------------------------------------------------------------
\section{Dal rig all'animazione: gameplay e cinematiche}
\label{p4:sec:animazioni}

Arrivati al personaggio modellato, nel corso si comincia a \textbf{riggarlo}:
prima si mettono le \textbf{ossa}, poi i \textbf{controlli}, si fanno le
\textbf{blend shape}, fino ad arrivare a un personaggio funzionante. Quando
il personaggio è riggato\footnote{Nella trascrizione ``quando il personaggio
è arricchito'': con ogni probabilità ``riggato''.}, entra in gioco
l'\textbf{animatore}, che si occupa di due tipi di animazione:
\begin{itemize}
  \item le \textbf{animazioni gameplay} (\emph{in-game animation});
  \item le \textbf{animazioni cinematiche}.
\end{itemize}
Il corso fa anche un po' di gameplay animation, soprattutto al primo anno, ma
si concentra sulle \textbf{cinematiche}, perché sono più complesse: ``se sai
fare animazioni cinematiche, quasi al 100\% sai fare animazioni gameplay''.

\subsection{Le animazioni gameplay: i loop}

Le animazioni gameplay sono tendenzialmente dei \textbf{loop}. C'è tutta la
\textbf{locomozione}: idle, camminata, giro a destra, giro a sinistra, corsa,
le transizioni da idle a camminata, da idle a corsa, da camminata a corsa, lo
stop da corsa a fermo, da corsa a camminata. Più il videogioco è complesso,
maggiore è il numero di animazioni richieste anche solo per la locomozione,
per gli attacchi e così via.

L'esempio è \emph{The Last of Us}, gioco tripla A con motion capture: lo
stesso personaggio ha \textbf{decine di idle} diverse. Sono tutti loop,
programmati in modo che, magari in maniera casuale o a rotazione, dopo otto
secondi di personaggio fermo parta un idle, poi un altro. In alcuni casi due
animazioni sono \textbf{sovrapposte}: un idle di base e, sopra, una variante
(si aggiusta i capelli, gioca con un oggetto che ha in mano).

Anche gli idle sono \textbf{specifici del personaggio}. C'è una ragazza
(``se avete giocato, no spoiler'')\footnote{Con ogni probabilità Abby, in
\emph{The Last of Us Part II} (2020).}, un personaggio tosto, che ha degli
idle più forti, mascolini: si scrocchia le mani, si aggiusta la spalla, si
pulisce il naso, ``così come farebbe un pugile''. Anche solo con un idle o una
camminata si racconta il personaggio.

% ---------------------------------------------------------------------
\section{Il programma dei due anni}
\label{p4:sec:programma}

Il relatore descrive il programma per far capire che cosa gli studenti di
character animation possono offrire agli altri (tabella~\ref{p4:tab:programma}).

\subsection{Il primo anno}

Oltre alle basi di modellazione, il primo anno ha due filoni.

\paragraph{Animazione.} Si parte dalla \textbf{pallina che rimbalza}, e se ne
fanno vari tipi: palle con pesi diversi, la palla che fa una corsa a ostacoli,
la palla con la coda, la palla con personalità. Poi i \textbf{cicli di
camminata}: prima solo una palla con i piedi, poi tutto il corpo senza le
braccia, poi il corpo intero, poi la camminata con personalità, la corsa. Poi
idle, salto, attacco singolo, attacco doppio (combo, melee), ``dead
cycle''\footnote{Così nella trascrizione: forse l'animazione di morte
(\emph{death}), forse un altro ciclo.}. Tutte queste animazioni si mettono
insieme nell'ultimo modulo del primo anno, quello di \textbf{Unreal Engine}:
con una \emph{state machine} si impara a muovere il proprio personaggio in
Unreal con le proprie animazioni. Per questo gli studenti sono incentivati a
fare tutte le animazioni con lo stesso personaggio. Si comincia anche un po'
di \textbf{meccanica del corpo}, un po' di combattimento.

\paragraph{Rigging.} Si parte da esercizi semplici, come riggare uno
spazzolino o una ruota panoramica, e si arriva a riggare il corpo, prima in
maniera semplice, poi complessa, con l'aggiunta di \textbf{blend shape
correttive} (per quando si piega il gomito, per quando si alza il braccio).

\subsection{Il secondo anno}

\paragraph{Animazione.} Ci si concentra sulla parte \textbf{cinematica}:
animazione di recitazione (\emph{acting}), meccanica del corpo e combattimenti
con più personaggi, animazione di \textbf{creature} (quadrupedi, draghi
ecc.), e \textbf{motion capture e retargeting}, con un workshop fatto anche
dal vivo.

\paragraph{Rigging.} Si passa al rigging \textbf{facciale} e delle
\textbf{creature}, e alla programmazione in \textbf{Python} per
l'automazione, sempre di rig, e per i \emph{pipeline tools}.

\begin{table}[htbp]
\centering
\small
\begin{tabular}{@{}lp{5.2cm}p{5.2cm}@{}}
\toprule
 & \textbf{Animazione} & \textbf{Rigging} \\
\midrule
\textbf{1° anno} & pallina che rimbalza (pesi, ostacoli, coda,
personalità); cicli di camminata e corsa; idle, salto, attacchi, combo;
meccanica del corpo; tutto in Unreal con una state machine &
spazzolino, ruota panoramica; corpo semplice, poi complesso con blend shape
correttive \\
\addlinespace
\textbf{2° anno} & acting; meccanica del corpo e combattimenti a più
personaggi; creature; motion capture e retargeting (workshop dal vivo) &
facciale; creature; Python per l'automazione dei rig e i pipeline tools \\
\bottomrule
\end{tabular}
\caption{Il programma del corso di Character Animation come lo riassume il
relatore. Si aggiungono i moduli propedeutici di modellazione, Unity e
Unreal.}
\label{p4:tab:programma}
\end{table}

\subsection{Quando si collabora}

La collaborazione fra gli altri corsi e Character Animation avviene di solito
al \textbf{secondo anno di Character Animation}: a volte sono ragazzi di
secondo con ragazzi di secondo, altre volte ragazzi del secondo anno di
character con ragazzi del terzo di game design e del secondo di 3D art, ``un
po' mischiati''.

% ---------------------------------------------------------------------
\section{Entrare nell'industria}
\label{p4:sec:industria}

Nel corso ci sono delle slide che si discutono in classe su come entrare
nell'industria da animatore o da rigger. Il relatore ne mostra alcune, senza
entrare troppo nel dettaglio, perché ``ci stanno delle cose che si applicano
un po' a tutti'' (e promette di condividerle). Le linee guida riguardano:
muovere i primi passi, creare un portfolio, scrivere una \emph{cover letter}
e un \emph{resume} (il curriculum), affrontare un \textbf{test} se viene
richiesto (sezione~\ref{p4:sec:test}), creare il proprio \textbf{network}.
Valgono per tutte le discipline.

\subsection{Il portfolio online}

È importante creare e \textbf{mantenere aggiornato} il proprio portfolio
online. Chi farà 3D userà probabilmente \textbf{ArtStation}; chi fa cose che
devono essere viste in video, \textbf{YouTube}; chi fa programmazione magari
mette delle cose su \textbf{GitHub}, pubblico o privato.

\subsection{Il network}

Il network di contatti non è solo la quantità di persone che si conoscono dal
vivo, con cui si è studiato. Essendo un'industria che \textbf{si muove molto
in remoto}, è importante anche:
\begin{itemize}
  \item aggiungere tanti contatti su Instagram, Facebook, LinkedIn;
  \item non vergognarsi di \textbf{chiedere pareri} sul proprio lavoro;
  \item \textbf{dare feedback} agli altri sul loro lavoro, nella maniera più
        educata, proattiva e costruttiva possibile.
\end{itemize}
Tutte queste persone, potenzialmente più facilmente di un recruiter o di un
supervisore, stanno lavorando (o un domani lavoreranno) a qualcosa: potranno
fare il vostro nome, raccomandarvi o darvi delle dritte.

Oltre alla cura online, bisogna partecipare alle \textbf{fiere di settore},
alle game jam, agli eventi di game dev, conoscere più realtà possibili,
soprattutto le piccole (``che ne so, la fiera dello sviluppo del videogioco
di Padova: non esiste, però\dots''): aiutano a entrare in contatto con altre
persone e a essere ascoltati, in una dimensione umana. Per questo è utile,
anche per chi è contrario ai social nella vita privata, avere dei
\textbf{profili social professionali} (Facebook, Instagram ecc.; ``per ora
TikTok non ha preso, da questo punto di vista; LinkedIn sì''), per fare un po'
di self-promotion, seguire le persone, restare aggiornati sull'industria.

\subsection{Il curriculum}

Avere sempre un curriculum \textbf{sia in inglese sia in italiano}, ben
aggiornato. Per un principiante il curriculum deve stare in \textbf{una
pagina}, e da quella pagina si devono capire le esperienze formative e quelle
professionali. Per un principiante sono esperienze professionali anche
\textbf{progetti come quelli di gruppo} che si faranno in AIV, in cui si crea
un videogioco senza che nessuno paghi, ``in maniera speculativa'', e tutte le
esperienze assimilabili al settore.

Non è indispensabile né utile mettere altre esperienze. Se il curriculum di
uno studente dice ``ho fatto per tre anni l'agente immobiliare, per due anni
ho lavorato all'Eurospin come cassiere''\footnote{Nell'esempio il relatore usa
il nome di uno dei partecipanti; l'esperienza è inventata.}, non serve: meglio
un solo paragrafetto, ``esperienze precedenti al di fuori dell'industria'',
con ``ho fatto questo, questo e questo per $x$ anni'', giusto per far capire
che si è persone serie, che hanno già fatto altri lavori.

\subsection{La cover letter}

Soprattutto quando ci si candida a studi che interessano, una \textbf{cover
letter} personalizzata, mandata per mail, può aiutare. È una lettera di
intenti, di interesse: ``Caro recruiter di One-O-One Games o di Milestone, da
quando sono piccolo gioco ai giochi MotoGP che sviluppate in Milestone'', o
``gioco al \emph{Suicidio di Rachel Foster}''; ``vi seguo da molto, ho
iniziato in questo settore anche sperando di collaborare con realtà come la
vostra; in particolare credo di essere un buon candidato perché le mie
competenze XYZ possono tornare utili al dipartimento XYZ''. Bisogna far capire
di aver \textbf{fatto i compiti a casa}.

% ---------------------------------------------------------------------
\section{Offerte di lavoro, stipendi, costo della vita}
\label{p4:sec:stipendi}

\subsection{Il foglio delle offerte}

Ci sono moltissimi siti su cui si postano offerte di lavoro, e c'è un
\textbf{Google Sheet} tenuto dalla community che ``non tanti conoscono, ma
secondo me è oro''\footnote{Il relatore ne dà il link in chat, ma la
trascrizione non riporta il nome del foglio.}. Contiene:
\begin{itemize}
  \item una pagina di \textbf{openings}: animatore, modellatore, surfacing,
        software engineer, producer, CFX (\emph{creature effects artist}),
        con l'indicazione se il lavoro è sul posto, remoto o ibrido e di
        quando è stato postato. Da agosto a quel momento le offerte sono
        \textbf{più di 3.000};
  \item una pagina \textbf{Studios}, con la lista degli studi di
        videogiochi, serie TV, cinema: più di 6.000 (6.600 in quel momento);
  \item una pagina \textbf{Internship}, con gli studi che cercano stagisti:
        140;
  \item le posizioni \textbf{entry level} e \emph{postgrad}, per chi si è
        appena diplomato; ma il relatore manderebbe sempre la candidatura
        anche se la posizione è per un \emph{mid}, ``perché comunque fai
        vedere la tua reel, magari piace'';
  \item una pagina \textbf{Education}: se si vuole fare un corso di FX, si
        cerca FX e si trova per esempio ArtFX\footnote{Scuola francese di
        effetti visivi, animazione e videogiochi.}, ``per l'FX una delle più
        fighe''; ci sono molte scuole e siti online dove fare workshop e
        proseguire la formazione dopo AIV, circa 150 voci;
  \item moltissimi siti di interesse: risorse, articoli, materiale di studio,
        texture, altre job board.
\end{itemize}

\subsection{Il foglio degli stipendi}

Insieme a questo c'è un secondo foglio, sugli stipendi\footnote{Nella
trascrizione ``Solti, Salary, Animator'': il nome non è ricostruibile con
sicurezza. È comunque un foglio compilato in forma anonima dalla community.}:
``un po' una manna dal cielo'', perché uno aiuta a trovare lavoro e l'altro a
farsi un'idea dei lavori che si trovano. Lì le persone inseriscono
\textbf{anonimamente} quanto guadagnano, per quale ruolo, con quale livello di
seniority, in che città; alcuni dicono anche l'azienda, il genere, se fanno
parte di un sindacato (\emph{union}).

Gli esempi letti a schermo, per offerte di Ubisoft:
\begin{itemize}
  \item un concept artist a \textbf{Bucarest}, tre anni d'esperienza,
        \emph{regular} (junior-mid): 29.000 lordi, che risultano in 17.000
        netti;
  \item un cinematic animator \textbf{junior in Italia}: circa 2.000 al
        mese;
  \item un gameplay animator junior in Italia: 28.000; e altri junior a
        24.000;
  \item un \emph{mid} con quattro anni d'esperienza: 27.500; un altro mid:
        27.000. In sintesi, Ubisoft Italia paga gli animatori fra i 24 e i
        28 mila;
  \item in \textbf{Svezia}, sempre Ubisoft: 38.000 corone svedesi, che il
        relatore traduce in 3.375 euro al mese, con cinque anni
        d'esperienza\footnote{Il relatore dice che le corone svedesi ``se non
        ricordo male sono abbastanza simili agli euro'': non è così, un euro
        vale circa 11 corone. La cifra di 3.375 euro al mese che legge è però
        coerente con 38.000 corone \emph{mensili} al cambio reale.}. Lì poi
        bisogna fare il conto del costo della vita.
\end{itemize}

Fra i junior italiani si vede una differenza di genere: nel caso specifico,
dice il relatore, probabilmente non ha pesato tanto il genere quanto la
\textbf{negoziazione} del contratto (sezione~\ref{p4:sec:genere}). Un sito del
genere fa capire che c'è chi prende 24 e chi 28. Quindi, se arriva un'offerta
da 24, invece di dire ``sì, sì, sono contentissimo'', si può rispondere:
``Grazie dell'offerta, la sto valutando molto seriamente. Mi piacerebbe
lavorare con voi. Al momento le mie aspettative salariali sarebbero più
vicine ai 30.000, ma possiamo sicuramente trovare un accordo. Qual è il limite
che vi siete imposti per questa posizione? Si può fare qualcosa di meglio?''.
Se rispondono ``no, assolutamente solo 24'', potrebbe voler dire che ci sono
persone con una reel un po' migliore disposte a lavorare per quella cifra: si
accetta, ``allora dopo ci provate''.

\subsection{Il costo della vita: Numbeo}

Il terzo sito è \textbf{Numbeo}, per confrontare il costo della vita, anche
fra due città. Esempio: si vive a Roma e si riceve un'offerta da Ubisoft in
Svezia, a \textbf{Malmö}. Il confronto letto a schermo\footnote{Le cifre di
Numbeo variano nel tempo e sono riportate come il relatore le legge; la
trascrizione le scrive in modo confuso (``17,35'', ``20,678'', ``7,80'').}:
\begin{itemize}
  \item a Malmö la vita costa in generale il 7,5\% in più di Roma, ma
        includendo l'affitto il costo diventa più basso: a Roma l'affitto
        pesa molto di più;
  \item la spesa costa di più a Malmö; il potere d'acquisto, però, a Malmö è
        il \textbf{56\% più alto};
  \item stipendio medio netto: circa 1.735 euro a Roma, 2.068 a Malmö;
  \item affitto di un monolocale fuori dal centro: 780 a Roma, 720 a Malmö;
  \item per mantenere a Malmö lo stile di vita che a Roma si ha con 4.800
        euro, ne servono 4.735: più o meno il costo della vita si pareggia.
\end{itemize}
La cosa interessante è che alcune voci sono marginali e ci si può
organizzare: il cappuccino a Malmö costa molto di più, ``porto la macchinetta
dall'Italia e lì mi risparmio''; in un ristorante non troppo costoso si paga
persino meno, quindi ``magari mi concedo una roba in più''\footnote{Poco
prima, però, il relatore aveva detto che ``i ristoranti costano molto di
più'': dalla trascrizione non si capisce se a Roma o a Malmö.}; il cellulare
costa di più, quindi si valuta un servizio di eSIM\footnote{Nella trascrizione
``Iralo, Sali'': con ogni probabilità Airalo e Saily.}; se si hanno bambini
conviene, perché costa molto meno.

Usando questi tre siti insieme ``esce fuori qualcosa di molto interessante'':
vanno consultati regolarmente. E sull'azienda in questione ci si informa
anche su \textbf{Indeed} e \textbf{Glassdoor}, dove ci sono recensioni su
quanto paga, quanto è seria e così via.

% ---------------------------------------------------------------------
\section{Età, provenienza, pseudonimi e portfolio}
\label{p4:sec:eta}

Il relatore mostra le indicazioni di portfolio, più legate ad animazione e
rig, ma alcune valgono per tutti: per esempio, il link deve essere
\textbf{aperto a tutti, senza password}.

\subsection{Quanto conta l'età? Nella sua esperienza, zero}

Federico chiede in chat quanto valore si dia all'età. ``Nella mia esperienza,
zero''. Per spiegare perché, il relatore racconta la propria biografia.

Ha iniziato a lavorare molto presto: viene da una famiglia non benestante, la
madre era ragazza madre con due figli, e lui ha cominciato a lavorare
\textbf{a 14 anni}. Studiava grafica pubblicitaria e faceva già loghi,
locandine, piccoli manifesti locali; poi è passato al web design.
``Ho anticipato tanto le tappe'': al quinto anno delle superiori si è
ritirato da scuola, ma ha comunque fatto il diploma nella stessa scuola, con
i suoi compagni di classe, come candidato privatista esterno, e si è iscritto
subito a un corso di 3D ``tipo AIV'' all'\textbf{Istituto Quasar} di Roma
(``era il 2008, credo''), proprio per anticipare ancora di più. Si è trovato
così molto giovane nell'industria.

Dopo il Quasar si è iscritto ad \textbf{Animation Mentor}. Nessuna delle due
scuole se la poteva permettere: ha preso un prestito di cinque anni con
Findomestic e uno di altri tre con Agos, ``pagando tipo il 13\% di
interessi, tassi quasi da strozzinaggio''. Nel frattempo lavorava: un anno da
\textbf{Mondo Convenienza} come venditore, usando però le competenze 3D per
il design degli ambienti; un anno come \textbf{3D generalist} e grafico
pubblicitario in uno studio fotografico, con clienti che chiedevano piccoli
lavori 3D (modellazione di lampade da esterno, modelli per manuali di
istruzione ``tipo quelli di IKEA''). Con quello che guadagnava pagava i
prestiti e campava. Ha anche avuto figli molto presto: una figlia di 16 anni,
avuta a 19, e un figlio di 13.

Ad Animation Mentor si è scontrato con una cosa. Era abituato a essere il
primo della classe, sempre in anticipo; lì c'erano persone di tutte le età,
anche di 32 anni quando lui ne aveva 19-20, con già 7, 8, 10 anni di
esperienza di animazione, che facevano cose meglio di lui. Il relatore mostra
il profilo LinkedIn di uno di loro, ``uno degli animatori più bravi che io ho
conosciuto nella mia carriera'' (``io LinkedIn lo uso sia per stalkerare le
persone, sia per farmi un'idea''): aveva fatto Animation Mentor con lui, o
poco prima, intorno al 2011, ed era già abbondantemente sopra i 30 (oggi è
``intorno ai 50''). Aveva cominciato come 3D animator nel \textbf{1998}, negli
Emirati Arabi, a Dubai, da dove proveniva: arrivato ad Animation Mentor aveva
già circa 12 anni d'esperienza. Oltre ad Animation Mentor ha seguito corsi ad
AnimSquad, ``una scuola che non esiste più''. Nel \textbf{2015} viene assunto
alla \textbf{Walt Disney}: per il suo \emph{big break} ci ha messo 17 anni,
dal 1998 al 2015, ed è arrivato in Disney che aveva forse 40 anni.

``La cosa bella del nostro settore'' è proprio questa: si manda una
\textbf{showreel}, e dalla showreel di solito non si sa di che genere sei,
quanti anni hai, se hai esperienze pregresse (o meglio, lo si può capire da
quello che c'è dentro). Si vedono solo le competenze. Quindi all'età non si
dà valore. Dipende magari dalle leggi dei singoli paesi: in Italia, per
esempio, una persona \textbf{under 35} dà diritto a degli sconti sul
contratto\footnote{Il relatore si riferisce agli incentivi all'assunzione di
giovani previsti dalla normativa italiana, la cui soglia d'età e la cui forma
sono cambiate più volte negli anni; il dettaglio non è verificabile dalla
trascrizione.}. Ma vale solo ``se non sei bravissimo'' o non hai un gran
portfolio: se allo studio servi, trovano il modo, al massimo si negozia sul
contratto.

\subsection{La provenienza e gli incentivi}

Altre cose, a volte, contano di più, come la \textbf{provenienza}. Un amico
del relatore, che ha cominciato da grande ad Animation Mentor, nel
\textbf{2013} aveva ottenuto una \textbf{internship in Pixar}: l'avevano
selezionato e gliel'avevano offerta, ma poteva farla solo trovando da sé il
modo di ottenere un visto americano, perché Pixar non sponsorizzava i visti
per gli stage. Non l'ha potuta fare.

Oggi questa cosa è molto comune soprattutto nei film e nelle serie, che
reggono molto sugli \textbf{incentivi statali, governativi, regionali}. Al
relatore è capitato anche fra regioni italiane: ``abbiamo preso un bando
della Regione Campania, quindi possiamo assumere solo persone con la
residenza in Campania''; o ti sposti, o non ti possono prendere, perché i
soldi con cui ti pagano vengono dalla Regione. Oppure uno studio di
Milano\footnote{Nella trascrizione ``202 Dogs Studio'': il nome non è
verificabile.} che assume chiedendo di andare a lavorare nel suo studio di
Fuerteventura, alle Canarie, perché ha preso un \emph{tax credit} delle
Canarie. (``Questo è: eravamo quattro amici al bar''.)

\subsection{Niente pseudonimi}

Per il portfolio online, il consiglio è \textbf{non usare pseudonimi o
alias}. L'unica eccezione è se si è già molto conosciuti online con lo
pseudonimo. Esempi:
\begin{itemize}
  \item l'amico islandese si chiama Sigurdur, ma tutti lo conoscono come
        ``Sighi'', perché il nome è difficile da pronunciare e scrivere;
  \item Daniele Angelozzi, scultore e istruttore ZBrush italiano, è da tutti
        conosciuto come ``Danko'': prima online si chiamava solo così, ora ha
        inserito il nome vero con Danko in mezzo;
  \item Francesco Lupo per tanto tempo online si chiamava ``Franchino'', e
        cercandolo lo si trova ancora con il vecchio nome.
\end{itemize}
Ma se si sta costruendo un'immagine pulita, meglio \textbf{nome e cognome}.
Immaginate che una studentessa raccomandi un amico al suo studio: ``c'è un mio
amico molto bravo, Matteo G.''. Prima ancora che lei condivida i contatti, il
supervisore cerca il nome online; non trova niente; ``basta, ce lo siamo
scordato, non esiste più''. Magari su Instagram lui si chiama
``Matteo\_Fragolino\_2009'': ``e basta, fine, non abbiamo più possibilità di
lavorare''. Bisogna anche tenere sempre aggiornate le informazioni su dove si
lavora, su cosa si fa, sulle cose imparate, sui nuovi corsi seguiti.

\subsection{Qualità, non quantità}

Nel portfolio puntare alla \textbf{qualità e non alla quantità}. Per
l'animazione, anche solo \textbf{tre scene animate} di solito bastano a
capire se una persona può fare bene questo lavoro; per il rigging, anche
\textbf{uno o due rig}. Mettere lavori che mostrino esattamente i propri
\textbf{punti di forza}; i punti di debolezza, in uno slancio di onestà,
menzionarli solo durante il colloquio. Se chiedono quali sono, si può dire:
``Sento di dover crescere ancora in questa cosa; la so fare, sono
indipendente, però è una cosa che voglio migliorare. Infatti sto studiando
quello in questo momento''. Questo perché la reel \textbf{non la si può
discutere} con chi la guarda: ``la guardano in muto, in due secondi, e ciao,
è finita''. E usare una \textbf{risoluzione adeguata}.

% ---------------------------------------------------------------------
\section{Genere e negoziazione}
\label{p4:sec:genere}

Una studentessa fa una ``domanda scomoda'': si fa molta differenza fra uomini
e donne? Lei su Instagram ha rimosso il nome e tenuto solo il cognome, per
vedere come cambiava la percezione, e forse per caso ha cominciato a trovare
molto più facilmente lavoro in un settore che le interessava,
l'horror\footnote{Ricostruzione della domanda, che nella trascrizione è
confusa (``ho cominciato a trovare un settore che mi interessava orrore molto
più facilmente'').}.

Nella sua ``personalissima esperienza'', risponde il relatore, no. Nel suo
studio numericamente ci sono sempre state più donne. Ma, essendo la misoginia
``una sorta di malattia sistemica'', non è impossibile che la differenza
esista: l'industria è internazionale, ci si confronta con studi di culture
diverse (gli Emirati Arabi, alcuni paesi balcanici) dove la differenza può
essere maggiore; e anche in Italia, assumendo per una produzione di tre o sei
mesi, ``magari ci può stare il datore di lavoro che fa un ragionamento più
misogino: quella magari ha i figli''. Nella sua esperienza non è così.

In animazione le donne sono tante, ``forse tante quante gli uomini, se non di
più''. Nell'industria si sono avvicinate molto negli ultimi dieci anni: per un
periodo era una cosa ``super nerd e super da smanettoni'', i software erano più
difficili e meno interessanti, e anche giocare ai videogiochi era visto
culturalmente in modo diverso per le donne. Ora crede che siano pari; nel
\textbf{rigging}, invece, ``è difficilissimo trovare una donna''. Quindi
nella sua esperienza non c'è stata discriminazione diretta; ma se tenere solo
il cognome funziona, perché no. Fermo restando che, se avesse contezza di una
discriminazione effettiva, la combatterebbe e cercherebbe di non lavorare in
quel posto: mettere solo il cognome può far arrivare un'offerta che altrimenti
non sarebbe arrivata, ma se la ragione per cui non era arrivata era l'essere
donna, ``probabilmente non ci vorresti lavorare, con quelle persone''. Diventa
un modo di filtrare ``senso comune e decenza''. La studentessa conferma che
lei si è trovata meglio così.

\subsection{Il tipo di progetti}

Può darsi, aggiunge il relatore, che sia più facile che alle donne facciano
fare cose per bambini, ma crede che dipenda molto dal \textbf{portfolio}: se è
già tendente all'horror, difficilmente propongono cose per bambini, ``perché
hanno paura che le fai un po' macabre, alla Tim Burton''; se invece il
portfolio è ``puccioso'', è più difficile che arrivino cose horror.

\subsection{La differenza salariale}

Differenze salariali fra uomini e donne il relatore ne ha viste spesso. Le
riconduce alla \textbf{negoziazione}, ``sempre un po' un male culturale,
potenzialmente riconducibile al patriarcato'': soprattutto dopo i primissimi
anni d'esperienza, all'uomo offrono 10 e lui chiede 12, a prescindere, anche
se 10 è più di quanto si aspettava (aspettativa di 8, offerta di 10, ``proviamo
con 12'', magari ne dà 11); le donne lo fanno un po' meno, forse perché
ricevono meno proposte. Alla fine può nascere una differenza.

Non è solo una questione di genere ma anche di \textbf{provenienza}. In
Canada gli europei occidentali venivano pagati più dei sudamericani e degli
indiani, perché questi volevano, tra virgolette, scappare dal loro paese e
andare a vivere altrove in maniera stabile; gli europei invece ragionavano
``ho degli studi fighi pure qua, campo bene pure qua; se mi fanno sorridere
vado in Canada, se no non è che sto sotto ai ponti''. Agli europei, racconta,
davano un soprannome volgare (``le puttane europee''), perché se tre mesi
dopo un altro studio offriva il 20\% in più ``facevi il ciaone'' e andavi
nell'altro studio. I canadesi, che magari stavano crescendo una famiglia là,
avevano bisogno di stabilità; indiani, sudamericani e altri davano almeno due
o tre anni di \emph{loyalty}, di fedeltà allo studio.

% ---------------------------------------------------------------------
\section{I test}
\label{p4:sec:test}

Delle slide resta da parlare del \textbf{test}, ``qualcosa che capita anche a
chi fa altro'' (le slide contengono anche una lista di eventi di animazione e
di videogiochi, e consigli per le review che il relatore salta). Il sentimento
dell'industria sui test è ``un po' ibrido'', perché è molto facile che
diventino una sorta di \textbf{scam}. Se uno studio chiede un test, è una
richiesta lecita, ma bisogna mettere in chiaro alcune cose \textbf{fin da
subito}:
\begin{itemize}
  \item lo studio deve dire se il test verrà \textbf{usato in produzione}. Se
        sì, va \textbf{pagato}, a prescindere dall'assunzione. Un accordo
        possibile: se passo il test sono assunto, e il test viene usato in
        produzione; se non lo passo, ma vi sembra abbastanza buono da
        chiedermi il file sorgente, me lo pagate, perché potreste usarlo in
        produzione senza neanche menzionarmi.
\end{itemize}
Come affrontarlo:
\begin{itemize}
  \item restare \textbf{aderenti alle specifiche} fornite;
  \item \textbf{stare nei tempi}: spesso il test non misura solo se si sa fare
        una cosa di una certa qualità, ma se la si sa fare \textbf{in tempi di
        produzione};
  \item fare \textbf{tante domande}, prima e durante: come vuoi che lo
        presenti, in che formato, dove te lo invio; per sembrare sempre sul
        pezzo;
  \item \textbf{lavorare subito in modo pulito}: niente nomi di file con
        bestemmie ``pensando che tanto non lo dovranno vedere'', perché magari
        poi chiedono il file e ci si dimentica di cambiarli; niente lavoro
        disordinato; una struttura di cartelle professionale.
\end{itemize}
Il resto delle slide ripete consigli già dati sul network.

% ---------------------------------------------------------------------
\section{Le domande finali}
\label{p4:sec:domande}

``Tutto ciò che io volevo dirvi in maniera orale ve l'ho detto'': il resto
dipende dalle domande, e la lezione ``può finire prima se voi non avete niente
da chiedere, o può durare a oltranza''.

\subsection{Che cosa si chiedono un narrative designer e un animatore?}

Edoardo, che fa narrative design (capitolo~\ref{p3:cap:narrative}), chiede che
richieste può ricevere da chi fa character animation e che richieste può
fargli. Il relatore fa l'esempio dello sceneggiatore televisivo con gli
animatori, il lead animator e il supervisor: c'è una \textbf{relazione
bidirezionale}.
\begin{itemize}
  \item L'animatore può dire quali sono i \textbf{limiti tecnici} del
        personaggio: ``il mio personaggio ha un rig che non gli permette di
        fare cose troppo movimentate''. Allora si evita di scrivere un
        personaggio grasso troppo movimentato e si punta sulla sua lentezza,
        sulla sua goffaggine.
  \item O viceversa: il personaggio, nonostante sia grosso, deve mostrare
        tanto dinamismo; allora il narrative designer chiede al rigger e
        all'animatore che il modello sia grosso, ma in modo da potersi
        permettere le capriole, perché è importante.
\end{itemize}
Il relatore immagina questa interazione soprattutto nello studio del
personaggio e degli ambienti. Esempio paradossale: un gioco inclusivo con un
personaggio in sedia a rotelle, e il narrative designer che scrive un
ambiente tipo la Terra di Mezzo, tutto montagne senza strade. Ma escono anche
cose più piccole. Il narrative designer scrive un livello in cui il
personaggio deve fare tante cose, fra cui una piccola ma indispensabile:
cogliere una mela da un albero. Mesi dopo arriva l'animatore; nel frattempo
sono state prese tante scelte, e il personaggio è troppo piccolo, non arriva
alla mela neanche saltando. Bisogna trovare una soluzione, per esempio una
roccia su cui salire, che però causa più lavoro al team degli asset e alla
programmazione. È l'interazione più comune, anche se i due dipartimenti
lavorano a una certa \textbf{distanza cronologica} e non sempre hanno una
collaborazione troppo stretta.

\subsection{Quanto conta il talento?}

Sara chiede quanto sia necessario essere talentuosi e quanto invece si
impari. Secondo il relatore:
\begin{itemize}
  \item per \textbf{lavorare}: 99\% fatica, 1\% talento;
  \item per lavorare ad \textbf{altissimi livelli}: 75 e 25;
  \item per lavorare ad altissimi livelli ed essere considerati ``un drago'',
        una leggenda nello studio e intorno: probabilmente 50 e 50.
\end{itemize}
``Se io voglio solo fare i videogiochi, non c'ho neanche bisogno di talento'':
tutto ciò che viene visto come talento si impara e si sviluppa, il gusto, la
capacità artistica, la capacità tecnica. Per lavorare ai giochi di Naughty
Dog probabilmente bisogna essere molto forti: o una buona quota di talento, o
tanti anni di esperienza cercando di migliorare sempre la produzione a cui si
lavora.

\subsection{Bisogna saper disegnare, modellare, programmare?}

Giulia chiede quanto sia fondamentale, per animazione e rigging, saper
disegnare e modellare. Il relatore aggiunge anche la programmazione:
\begin{itemize}
  \item \textbf{programmare} è solo un plus per l'animazione, mentre per il
        rigging oggi è \textbf{quasi indispensabile};
  \item \textbf{disegnare e modellare} non sono indispensabili per nessuna
        delle due, ma entrambe aiutano: disegnare aiuta di più l'animazione,
        modellare aiuta di più il rigging.
\end{itemize}
Chi volesse specializzarsi nel rigging dovrebbe modellare un pochino e
programmare un pochino; chi fa animazione, puntare più sul disegno. Però
l'animazione ``non è una tecnica, è una disciplina'', come il disegno: se già
si sa disegnare aiuta, ma se non si sa, i primi anni sono meglio spesi
imparando l'animazione. È pieno di animatori che, già al lavoro, seguono corsi
di disegno dal vivo ``per sciogliere un po' la mano'', perché poi servono
solo degli schizzi.

\subsection{3D art e character animation vanno a braccetto?}

Alessandro chiede se i corsi di 3D art e character animation vadano a
braccetto. Non sempre, ma a volte si parlano, soprattutto nei progetti
interdisciplinari\footnote{Nella trascrizione ``the biggest ones, big ones'',
e all'inizio ``gli ex biggest ones, ovvero questi progetti tesi in cui vi
mettete insieme con vari dipartimenti per tirare fuori il vostro gioco da
lanciare su Steam'': con ogni probabilità sono i Dev Projects
(sezione~\ref{p1:sec:devprojects}).}: spesso il character modellato in 3D art
viene riggato e animato dai ragazzi del secondo anno di character animation.

\subsection{Gli studenti di 3D possono imparare un po' di rigging e animazione?}

Infine Giulia chiede se gli studenti di 3D potranno imparare anche un po' di
rigging e animazione. Va chiesto ad \textbf{Andrea Morini}: potenzialmente
potrebbe mettere a disposizione le registrazioni degli anni passati, o far
partecipare ad alcune lezioni del corso, ma il relatore non può dare una
risposta concreta perché servirebbe un'organizzazione che ricade sulla
scuola. Lui sarebbe aperto ad avere studenti di altre classi; in alternativa
si può chiedere se è possibile fare un piccolo modulo. Per i corsi esterni,
gli si può chiedere in qualsiasi momento.

\subsection{I saluti}

``Mi sembra di capire che avete fame e volete andare a pranzo''. Il relatore
resta a disposizione, anche online e durante i corsi, e augura a tutti un buon
inizio d'anno: ``in bocca al lupo e buon proseguimento''.
